% sec_mechanism_evidence.tex — V8 compressed

\subsection{Bid Rotation and Bid Inflation}
\label{sec:bid_rotation}

\paragraph{Winner persistence.}
FL firms' winner HHI is 0.178 (14.3 unique winners), significantly
lower than non-FL always-losers' 0.303 (5.0 winners; $p < 0.001$).
FL firms co-participate with a wider range of winners, consistent
with cartels rotating designated winners in competitive markets
(Proposition~\ref{prop:market_selection}).

\paragraph{Repeated FL--winner pairs.}
Among 38,941 FL--winner pairs, 4,603 share $\geq 5$ tenders and
379 share $\geq 20$ (maximum: 177). Non-FL pairs average only
1.46 shared tenders with 494 pairs reaching $\geq 5$. The heavy
right tail indicates persistent cartel relationships.

\paragraph{Bid inflation.}
FL median bid-to-winner ratio is 1.85 (85\% above the winner) vs.\
1.43 for non-FL losers. Controlling for item and year FE, FL bids
are 15.4\% higher ($p < 0.001$;
Figure~\ref{fig:bid_inflation} in Appendix~\ref{app:robustness})---not
merely unlucky competitive bids but strategically elevated submissions
designed to lose.
